\section{The endpoint obstruction defines a two-dimensional heat boundary}
\label{sec:endpoint-sectorial-heat}

\paragraph{ESH1. The original theta density before its differential filter.}
Keep the original coordinate \(s=\tfrac12+\tfrac{iZ}{2}\), and
the full factor \(A(s)=\pi^{-s/2}\Gamma(s/2)\). The theta
formula LH15, with the exact substitution \(x=e^{4u}\), gives
\[
 A(s)\zeta(s)=\int_0^\infty\rho(u)\cos(Zu)\,du,\qquad
 \rho(u)=8e^u\psi(e^{4u})-4e^{-u},\qquad |\Im Z|<1.
                                                               \tag{ESH1}
\]
Indeed the two integral terms in LH15 become
\(8\int_0^\infty e^u\psi(e^{4u})\cos(Zu)du\).
Its rational endpoint contribution is
\(1/[s(s-1)]=-4/(1+Z^2)\), equal to the retained
\(-4e^{-u}\) integral on the stated strip. Both endpoint
terms are consequently present in ESH1. The density and the
subsequent theta filter use the original definitions of
\href{https://arxiv.org/abs/1801.05914v5}{Brad Rodgers and
Terence Tao, \emph{The de Bruijn--Newman constant is non-negative}},
equations \texttt{phidef}, \texttt{htdef}, \texttt{hoz} and
\texttt{sas}; the unfiltered sectorial extension is derived here.

The original unfiltered real-axis integral after inserting a
heat weight is
\[
 I_t(Z)=\int_0^\infty e^{tu^2}\rho(u)\cos(Zu)\,du.
                                                               \tag{ESH2}
\]
It converges for \(\Re t<0\), for every complex \(Z\),
locally uniformly together with all derivatives. The part
\(8e^u\psi(e^{4u})\) has faster than Gaussian decay, while
the endpoint part has Gaussian decay precisely in this left
half-plane. For positive real \(t\) and \(Z=0\), the latter
part diverges with one sign. Thus the original real-axis formula
does not define an entire-in-\(Z\) heat family at positive time.
An analytic continuation must be specified rather than calling
that same improper integral convergent.

\paragraph{ESH2. Exact lateral continuations of the endpoint.}
Set \(q=-t\), first \(q>0\). For any complex \(a\),
completing the square gives
\[
 F(q,a)=\int_0^\infty e^{-qu^2-au}du
 =\frac{\sqrt\pi}{2\sqrt q}
       e^{a^2/(4q)}\operatorname{erfc}\!\left(\frac a{2\sqrt q}\right).
                                                               \tag{ESH3}
\]
For real \(a\) this is a direct real integral substitution;
both sides are entire in \(a\) by Gaussian domination and
the entire power series of \(\operatorname{erfc}\), so the
identity holds for all \(a\). The right side extends analytically
to \(q\in\mathbb C\setminus(-\infty,0]\), with the square
root positive on the positive real axis. This specifies the branch.
Here
\(\operatorname{erfc}(z)=1-\frac2{\sqrt\pi}\int_0^z e^{-v^2}dv\);
its entire primitive immediately proves
\(\operatorname{erfc}(z)+\operatorname{erfc}(-z)=2\).
This is the convention of N. M. Temme,
\href{https://dlmf.nist.gov/7.2.E2}{NIST DLMF, equation 7.2.2};
the original equation TeX is retained with the source record.

The endpoint cosine integral is
\[
 B(q,Z)=\tfrac12\{F(q,1-iZ)+F(q,1+iZ)\}.                 \tag{ESH4}
\]
For positive real \(t\), define \(B_+(t,Z)\) and
\(B_-(t,Z)\) by approaching \(q=-t\) from the upper and
lower \(q\)-half-planes, respectively. The square roots are
\(+i\sqrt t\) and \(-i\sqrt t\). Thus ESH3 and the
displayed erfc identity give the exact difference
\[
 F_+(t,a)-F_-(t,a)
       =-i\sqrt{\frac\pi t}\,e^{-a^2/(4t)},
\quad
 B_+(t,Z)-B_-(t,Z)
       =-i\sqrt{\frac\pi t}\,
          e^{(Z^2-1)/(4t)}\cos\!\frac Z{2t}.              \tag{ESH5}
\]
This orientation uses the \(q\)-plane; reversing the order of
the two lateral values reverses the sign.

The rapidly decaying part
\[
 L_t(Z)=8\int_0^\infty e^{tu^2+u}\psi(e^{4u})\cos(Zu)du
\]
is jointly entire in \((t,Z)\). The two unfiltered continuations
at positive time are therefore
\[
 I_\pm(t,Z)=L_t(Z)-4B_\pm(t,Z),\qquad
 I_+-I_-=4i\sqrt{\frac\pi t}\,
                 e^{(Z^2-1)/(4t)}\cos\!\frac Z{2t}.
                                                               \tag{ESH6}
\]
The factor four and its sign come from the original endpoint
density in ESH1. Each branch satisfies
\(\partial_t I=-\partial_Z^2 I\): differentiate the integral
in \(\Re t<0\), then use analytic continuation on its
specified branch. This continuation assertion does not make
ESH2 converge at positive real time.

\paragraph{ESH3. The full theta filter and its exact kernel.}
Define, without suppressing its time-dependent terms,
\[
 \mathcal N_t=-\frac{1+(Z-2t\partial_Z)^2}{64}
 =-\frac{1+Z^2-2t-4tZ\partial_Z+4t^2\partial_Z^2}{64}.
                                                               \tag{ESH7}
\]
For a twice differentiable density with the decay of \(\rho\),
two integrations by parts in \(\Re t<0\) give
\[
 \int_0^\infty e^{tu^2}\rho''(u)\cos(Zu)du
  =-\rho'(0)-(Z-2t\partial_Z)^2 I_t(Z).                  \tag{ESH8}
\]
The first boundary at infinity vanishes; at zero it is
\(-\rho'(0)\). The second boundary is zero because the
derivative of \(e^{tu^2}\cos(Zu)\) vanishes at zero.
Differentiating that whole weight twice gives precisely
\((2t+4t^2u^2-Z^2)\cos(Zu)-4tZu\sin(Zu)\),
which verifies every term of ESH8.

The Poisson identity in LH14 yields
\(\psi'(1)=-\psi(1)/4-1/8\). Consequently
\[
 \rho'(0)=8\psi(1)+32\psi'(1)+4=0.                     \tag{ESH9}
\]
There has been no dropped endpoint boundary. Direct differentiation
of the entire density, with \(x=e^{4u}\), gives
\[
 \frac{\rho''(u)-\rho(u)}{64}
 =\sum_{n\ge1}
   \left(2\pi^2n^4e^{9u}-3\pi n^2e^{5u}\right)e^{-\pi n^2e^{4u}}
 =\Phi(u).                                               \tag{ESH10}
\]
The \(-4e^{-u}\) term is annihilated by \(\partial_u^2-1\),
but its boundary contribution has already been included in ESH9.
ESH8--10 prove, first in \(\Re t<0\) and then along each
continuation, the exact identity
\[
 \mathcal N_t I_+(t,Z)=\mathcal N_t I_-(t,Z)
        =H_t(Z)=\int_0^\infty e^{tu^2}\Phi(u)\cos(Zu)du.
                                                               \tag{ESH11}
\]
The integral for \(H_t\) is jointly entire by the decay of
\(\Phi\). This is a derived equality between the two specified
heat constructions, including their lost boundary directions.
The two density parts separately have nonzero boundary terms:
\[
 \mathcal N_t B=-\frac1{64},\qquad
 \mathcal N_t L_t=H_t-\frac1{16},\qquad
 \mathcal N_t(L_t-4B)=H_t.                              \tag{ESH17}
\]
For the first equation the endpoint density satisfies
\(\rho''=\rho\), \(\rho'(0)=-1\), so ESH8 gives it
directly. For the second, the theta-only density has derivative
\(-4\) at zero by ESH9 and the endpoint derivative \(+4\)
in the full density. Its interior filtered density is still
\(\Phi\). Substitution in ESH8 yields the displayed
\(-1/16\), which the endpoint contribution cancels exactly.
Thus annihilating the interior endpoint density does not authorize
omitting its integration boundary.

For each \(t\ne0\), the kernel of \(\mathcal N_t\) on the
entire functions of \(Z\) is exactly
\[
 \left\{
 e^{Z^2/(4t)}
 \left(c_+e^{iZ/(2t)}+c_-e^{-iZ/(2t)}\right):
          c_+,c_-\in\mathbb C\right\}.                   \tag{ESH12}
\]
To prove completeness of this list, write
\(f=e^{Z^2/(4t)}g\). Then
\((Z-2t\partial_Z)f=-2t e^{Z^2/(4t)}g'\), and the
kernel equation becomes \(g+4t^2g''=0\). Its two initial
values determine its entire solution by the coefficient recursion,
giving exactly the two displayed exponentials. ESH5 lies in this
space, so its annihilation in ESH11 is also verified directly.
The actual cosine source is even in \(Z\). On this original
even domain the kernel is one-dimensional: evenness in ESH12
requires \(c_+=c_-\), leaving
\(e^{Z^2/(4t)}\cos(Z/(2t))\). The sine direction belongs
to the explicitly larger, nonsymmetric entire-function domain.
Thus the actual lateral jump occupies exactly the symmetric
kernel direction, not two independently varied cosine parameters.
The time dependence of any such homogeneous difference which
also satisfies the heat equation is determined. Writing the
two coefficients in ESH12 as functions \(c_\pm(t)\), direct
differentiation gives
\[
 c_\pm'(t)+\left(\frac1{2t}-\frac1{4t^2}\right)c_\pm(t)=0,
 \qquad c_\pm(t)=C_\pm t^{-1/2}e^{-1/(4t)}.              \tag{ESH18}
\]
On a fixed simply connected time domain avoiding zero the chosen
square root makes this formula single-valued. The two independent
space exponentials force both scalar equations, and integrating
them gives exactly the stated constants. In original \(s\)
coordinates the modes are \(t^{-1/2}e^{-s^2/t}\) and
\(t^{-1/2}e^{-(s-1)^2/t}\). The even source requires equal
coefficients, recovering the time dependence of ESH6 and ESH14.

\paragraph{ESH4. Both original-zeta receivers and the endpoint modes.}
Return to the original coordinate \(Z=-2i(s-1/2)\) and keep
\[
 f_\pm(t,s)=\frac{I_\pm(t,-2i(s-1/2))}
                   {\pi^{-s/2}\Gamma(s/2)},\qquad
 z_t(s)=\frac{16H_t(-2i(s-1/2))}
                 {s(s-1)\pi^{-s/2}\Gamma(s/2)}.           \tag{ESH13}
\]
The first family is defined on the sectorial branches just proved;
the second is CF1's returned meromorphic family. Both recover the
original \(\zeta(s)\) at time zero on \(0<\Re s<1\):
for the first, approach zero through \(t<0\) and dominate the
endpoint integral by \(e^{-(1-|\Im Z|)u}\); for the second,
use ESH1 and \(16\mathcal N_0=s(s-1)\).
Their domains at time zero and their later poles have not been
identified merely from their common initial values.

The actual two endpoint modes in the original coordinate are
\(e^{-s^2/t}\) and \(e^{-(s-1)^2/t}\). Specifically,
\[
 f_+(t,s)-f_-(t,s)
  =2i\sqrt{\frac\pi t}\,
       \frac{\pi^{s/2}}{\Gamma(s/2)}
          \left[e^{-s^2/t}+e^{-(s-1)^2/t}\right].          \tag{ESH14}
\]
This follows by expanding the cosine in ESH6 and substituting
the stated coordinate, with no rescaling of time.

The complete differential receiving map is
\[
 z_t(s)=\frac{
  [s(s-1)+t/2]A(s)f_\pm(t,s)
   +\frac{t(2s-1)}2\partial_s[A(s)f_\pm(t,s)]
   +\frac{t^2}4\partial_s^2[A(s)f_\pm(t,s)]
 }{s(s-1)A(s)}.                                         \tag{ESH15}
\]
For verification, \(Z-2t\partial_Z=-i[(2s-1)+t\partial_s]\);
squaring retains the extra derivative term \(2t\).
Multiplication by 16 in ESH7 then gives exactly the numerator
of ESH15. This proves the equality away from the displayed
exceptional factors and hence meromorphically, with the original
factor germs retained. The first derivative contributes
\(A'f+Af'\), and the second contributes
\(A''f+2A'f'+Af''\), with
\[
 A'/A=-\tfrac12\log\pi+\tfrac12\psi_{\Gamma}(s/2),\qquad
 A''/A=\tfrac14\psi_{\Gamma}'(s/2)
       +[-\tfrac12\log\pi+\tfrac12\psi_{\Gamma}(s/2)]^2.
                                                               \tag{ESH16}
\]
Here \(\psi_{\Gamma}\) denotes the Gamma logarithmic derivative,
distinct from the theta sum \(\psi(x)\). Every derivative of
the original multiplier is therefore specified.

At fixed \(t\ne0\) on these branches, \(I_\pm\) is entire
in \(s\), as is \(1/A(s)\), so \(f_\pm\) is entire. In
contrast, CF27 proves that \(z_t\) has a nonzero pole at one
for real \(t\). ESH15 is the exact map between these different
analytic domains; its denominator and derivatives cannot be
deleted. The recovery of \(\zeta\)'s pole as \(t\to0\) for
the sectorial family is not locally uniform at that pole, since
a locally uniform limit of holomorphic functions would be
holomorphic there. This identifies the precise endpoint domain
change, not a failure of the proved convergence in the open strip.

Thus the obstruction ESH2 defines both a two-branch continuation
and the explicit two-dimensional entire kernel ESH12, whose
intersection with the actual even source is one-dimensional.
The map to the theta heat family annihilates exactly these
homogeneous directions on the respective stated domains at each
nonzero time. Keeping a chosen continued
source and its endpoint coordinates preserves them; the
subsequent map to the original \(\zeta\) is ESH13--16 with
all factors retained. No conclusion about RH follows merely
from forgetting or retaining this computed kernel.
